% Preámbulo modular para presentaciones Beamer (Modelación Estadística)
% Diseño y paleta institucional del Tecnológico de Monterrey

\documentclass[aspectratio=169,xcolor={dvipsnames,table}]{beamer}

\usepackage[utf8]{inputenc}
\usepackage[spanish,mexico]{babel}
\usepackage{amsmath,amssymb,amsfonts,mathtools}
\usepackage{graphicx}
\usepackage{listings}
\usepackage{booktabs}
\usepackage{tikz}
\usetikzlibrary{matrix,arrows,decorations.pathmorphing,shapes.geometric,positioning}

% Paleta Institucional Tecnológico de Monterrey
\definecolor{TecAzul}{HTML}{0039A6}       % Azul TEC: color primario institucional
\definecolor{TecAzulOscuro}{HTML}{1A2E51} % Azul marino oscuro
\definecolor{TecRojo}{HTML}{EC2661}       % Rojo de acento (paleta miTec)
\definecolor{TecGris}{HTML}{646464}       % Gris neutro para comentarios y subtítulos
\definecolor{TecFondoBloque}{HTML}{F4F6F9}% Fondo suave para bloques

% Configuración de Tema Beamer
\usetheme{Boadilla}
\usecolortheme[named=TecAzulOscuro]{structure}
\setbeamercolor{alerted text}{fg=TecRojo}
\setbeamercolor{palette primary}{bg=TecAzulOscuro,fg=white}
\setbeamercolor{palette secondary}{bg=TecAzul,fg=white}
\setbeamercolor{palette tertiary}{bg=TecAzulOscuro!80!black,fg=white}
\setbeamercolor{title}{fg=TecAzulOscuro,bg=TecFondoBloque}
\setbeamercolor{frametitle}{fg=TecAzulOscuro,bg=TecFondoBloque}

% Bloques personalizados elegantes
\setbeamercolor{block title}{bg=TecAzulOscuro,fg=white}
\setbeamercolor{block body}{bg=TecFondoBloque,fg=black}
\setbeamercolor{block title alerted}{bg=TecRojo,fg=white}
\setbeamercolor{block body alerted}{bg=TecRojo!8,fg=black}
\setbeamercolor{block title example}{bg=TecAzul,fg=white}
\setbeamercolor{block body example}{bg=TecAzul!8,fg=black}

% Redondear bordes y añadir sombra sutil
\setbeamertemplate{blocks}[rounded][shadow=true]
\setbeamertemplate{navigation symbols}{}

% Pie de página limpio con número de diapositiva
\setbeamertemplate{footline}{
  \leavevmode%
  \hbox{%
  \begin{beamercolorbox}[wd=.7\paperwidth,ht=2.25ex,dp=1ex,left]{author in head/foot}%
    \hspace*{2ex}\insertshortauthor~~(\insertshortinstitute) \hfill \insertshorttitle
  \end{beamercolorbox}%
  \begin{beamercolorbox}[wd=.3\paperwidth,ht=2.25ex,dp=1ex,right]{date in head/foot}%
    \insertshortdate{}\hspace*{2em}
    \insertframenumber{} / \inserttotalframenumber\hspace*{2ex}
  \end{beamercolorbox}}%
  \vskip0pt%
}

% Configuración de Listings de Python para visualización pedagógica de código
\lstset{
    language=Python,
    basicstyle=\ttfamily\footnotesize,
    keywordstyle=\color{TecAzulOscuro}\bfseries,
    stringstyle=\color{TecRojo},
    commentstyle=\color{TecGris}\itshape,
    showstringspaces=false,
    breaklines=true,
    frame=lines,
    rulecolor=\color{TecAzulOscuro!30},
    backgroundcolor=\color{TecFondoBloque},
    aboveskip=0.5em,
    belowskip=0.5em,
    literate={á}{{\'a}}1 {é}{{\'e}}1 {í}{{\'i}}1 {ó}{{\'o}}1 {ú}{{\'u}}1
             {Á}{{\'A}}1 {É}{{\'E}}1 {Í}{{\'I}}1 {Ó}{{\'O}}1 {Ú}{{\'U}}1
             {ñ}{{\~n}}1 {Ñ}{{\~N}}1 {¿}{{?`}}1 {¡}{{!`}}1
}

% Metadatos institucionales por defecto
\title[Modelación Estadística]{Modelación Estadística para Ciencia de Datos e Ingeniería}
\author[J. Castillo Colmenares]{Juliho David Castillo Colmenares}
\institute[Tec de Monterrey]{Tecnológico de Monterrey \\ Escuela de Ingeniería y Ciencias}
\date{\today}

% Comandos y macros estadísticas compartidas para las presentaciones Beamer (Metropolis)

% Conjuntos numéricos básicos
\newcommand{\R}{\mathbb{R}}
\newcommand{\N}{\mathbb{N}}
\newcommand{\Q}{\mathbb{Q}}
\newcommand{\Z}{\mathbb{Z}}
\newcommand{\set}[1]{\left\{ #1 \right\}}
\newcommand{\abs}[1]{\left|#1\right|}
\newcommand{\norm}[1]{\left\| #1 \right\|}

% Comandos estadísticos y probabilísticos del curso
\newcommand{\Var}{\operatorname{Var}}
\newcommand{\cov}{\operatorname{Cov}}
\newcommand{\corr}{\rho}
\newcommand{\comb}[2]{\begin{pmatrix} #1 \\ #2 \end{pmatrix}}
\newcommand{\s}{\sigma}
\newcommand{\particion}{\mathfrak{P}}
\newcommand{\card}[1]{\#\left( #1 \right)}
\newcommand{\seq}[3]{\set{#1}_{#2}^{#3}}
\newcommand{\E}{\mathbb{E}}
\newcommand{\Prob}{\mathbb{P}}

% Entornos pedagógicos y cajas en español académico natural
\newenvironment{discusion}{%
  \begin{alertblock}{Pregunta de Discusión}%
}{%
  \end{alertblock}%
}

\newenvironment{reflexion}{%
  \begin{alertblock}{Reflexión para el Aula}%
}{%
  \end{alertblock}%
}

\newenvironment{paradoja}{%
  \begin{alertblock}{Paradoja de Exploración}%
}{%
  \end{alertblock}%
}

\newenvironment{motivacion}{%
  \begin{exampleblock}{Motivación: Ciencia de Datos en la Práctica}%
}{%
  \end{exampleblock}%
}

\newenvironment{retoclase}{%
  \begin{block}{Reto Rápido en Equipo}%
}{%
  \end{block}%
}


\title[Probabilidad Condicional e Independencia]{Sección 02.04: Probabilidad Condicional e Independencia}
\subtitle{Actualización de Incertidumbre, Regla de la Multiplicación y Paradoja de Monty Hall}
\author[J. Castillo Colmenares]{Juliho Castillo Colmenares}
\institute{Tecnológico de Monterrey}
\date{}

\begin{document}

% Slide 1: Portada
\begin{frame}[plain]
  \titlepage
\end{frame}

% Slide 2: Hoja de Ruta
\begin{frame}{Hoja de Ruta — Capítulo 02: Teoría de Probabilidad}
  ¿Dónde nos encontramos en el desarrollo modular de la teoría? \pause
  \begin{itemize}\itemsep=0.1em
    \item \textbf{02.01 Introducción a la probabilidad} \emph{(Completado)} \pause
    \item \textbf{02.02 Notación de conjuntos y particiones} \emph{(Completado)} \pause
    \item \textbf{02.03 Fundamentos y axiomas de probabilidad} \emph{(Completado)} \pause
    \item \color{TecRojo}\textbf{02.04 Probabilidad condicional e independencia \emph{(Hoy)}}\color{black} \pause
    \item \textbf{02.05 Teorema de Bayes} \pause
    \item \textbf{02.06 Muestreo aleatorio}
  \end{itemize}
  \vspace{0.4em}
  \begin{block}{Objetivo de la Sección 02.04}
    Formalizar el mecanismo matemático para actualizar la probabilidad de un evento ante la ocurrencia confirmada de otro, derivar la regla general de la cadena y distinguir rigurosamente entre independencia y mutua exclusión.
  \end{block}
\end{frame}

% Slide 3: Motivación
\begin{frame}{El Rol del Conocimiento Previo en Ciencia de Datos}
  \begin{motivacion}[¿Por qué necesitamos actualizar probabilidades?]
    Si sabemos que un correo electrónico contiene las palabras \texttt{verificación} y \texttt{urgente}, la probabilidad de que sea \emph{spam} no es la misma que para un correo cualquiera. El conocimiento parcial transforma radicalmente nuestra evaluación de incertidumbre.
  \end{motivacion}
  
  \vspace{0.3em}
  \pause
  \begin{itemize}\itemsep=0.1em
    \item \textbf{Inferencia Dinámica:} Los modelos de diagnóstico médico, detección de fraude o filtrado algorítmico operan evaluando eventos condicionados a evidencia observada. \pause
    \item \textbf{Reducción del Espacio:} Al confirmar que un evento $A$ ha ocurrido, el espacio muestral original $S$ se contrae, y $A$ se convierte en el nuevo universo de referencia.
  \end{itemize}
\end{frame}

% Slide 4: Definición Formal
\begin{frame}{Definición Formal de Probabilidad Condicional}
  \small Sean $A, B \subset S$ dos eventos tales que $P(A) > 0$. La \emph{probabilidad condicional} del evento $B$ dado que el evento $A$ ha ocurrido, denotada por $P(B \mid A)$, se define exactamente como: \pause

  \vspace{0.2em}
  \begin{block}{Definición (Probabilidad Condicional)}
    $$P(B \mid A) = \dfrac{P(A \cap B)}{P(A)}$$
  \end{block}

  \vspace{0.1em}
  \pause
  \begin{alertblock}{Interpretación Geométrica y Reescalamiento}
    Como sabemos con certeza que el resultado está dentro de $A$, los únicos resultados favorables para $B$ que aún pueden ocurrir son aquellos en la intersección $A \cap B$. Para mantener la normalización unitaria (Axioma 2), dividimos entre la probabilidad total de la nueva referencia $P(A)$.
  \end{alertblock}
\end{frame}

% Slide 5: Ejemplo Clásico Teórico
\begin{frame}{Ejemplo Clásico: Urna sin Reemplazo (1/2)}
  \small
  \begin{block}{Enunciado del Ejemplo Teórico}
    En un bote se colocan cinco canicas: tres blancas y dos negras ($N = 5$). Inmediatamente después se agita el bote. Determina la probabilidad de que en el segundo turno extraigamos una canica negra ($E_2'$), dado que en el primero obtuvimos una blanca ($E_1$).
  \end{block}

  \vspace{0.2em}
  \pause
  \begin{itemize}\itemsep=0.1em
    \item Sea $E_1$: \texttt{``Obtener canica blanca en 1er intento''} ($P(E_1) = 3/5$). \pause
    \item Sea $E_2'$: \texttt{``Obtener canica negra en 2o intento''}. \pause
    \item La probabilidad conjunta de sacar primero blanca y después negra es:
    $$P(E_1 \cap E_2') = \dfrac{3}{5} \times \dfrac{2}{4} = \dfrac{6}{20} = \dfrac{3}{10}$$
  \end{itemize}
\end{frame}

% Slide 6: Ejemplo Clásico Teórico (Continuación)
\begin{frame}{Ejemplo Clásico: Urna sin Reemplazo (2/2)}
  \small
  Aplicando directamente la definición formal de probabilidad condicional: \pause

  \vspace{0.2em}
  \begin{alertblock}{Cálculo Condicional Exacto}
    $$P(E_2' \mid E_1) = \dfrac{P(E_1 \cap E_2')}{P(E_1)} = \dfrac{3/10}{3/5} = \dfrac{1}{2} = 0.50$$
  \end{alertblock}

  \vspace{0.2em}
  \pause
  \begin{block}{Observación (Validez Axiomática)}
    La función $P(\cdot \mid A)$ satisface los tres axiomas de Kolmogorov sobre la subálgebra de eventos:
    $$P(A \mid A) = 1, \quad P(B \mid A) \geq 0, \quad P(B_1 \sqcup B_2 \mid A) = P(B_1 \mid A) + P(B_2 \mid A)$$
  \end{block}
\end{frame}

% Slide 7: Regla de la Multiplicación (o de la Cadena)
\begin{frame}{Teorema 2.4.1: Regla de la Multiplicación}
  \footnotesize
  Despejando la probabilidad conjunta de la definición condicional, obtenemos la herramienta fundamental para secuencias de experimentos: \pause

  \vspace{0.1em}
  \begin{block}{Regla de la Multiplicación para Dos Eventos}
    $$P(A \cap B) = P(A) P(B \mid A) = P(B) P(A \mid B)$$
  \end{block}

  \vspace{0.1em}
  \pause
  \begin{block}{Regla General de la Cadena para Tres o Más Eventos}
    Para cualesquiera tres eventos $A_1, A_2, A_3 \subset S$ con probabilidad positiva:
    $$P(A_1 \cap A_2 \cap A_3) = P(A_1) P(A_2 \mid A_1) P(A_3 \mid A_1 \cap A_2)$$
    En general, para $n$ eventos secuenciales:
    $$P(A_1 \cap \dots \cap A_n) = P(A_1) \prod_{k=2}^{n} P(A_k \mid A_1 \cap \dots \cap A_{k-1})$$
  \end{block}
\end{frame}

% Slide 8: Teorema de la Probabilidad Total
\begin{frame}{Teorema 2.4.2: Ley de la Probabilidad Total}
  \small Si $S = A_1 \sqcup A_2 \sqcup \dots \sqcup A_N$ es una partición disjunta del espacio muestral donde $P(A_i) > 0$ para todo $i$, entonces cualquier evento $B$ puede reconstruirse sumando sus contribuciones condicionales: \pause

  \vspace{0.2em}
  \begin{block}{Fórmula de la Probabilidad Total}
    $$P(B) = \sum_{i=1}^{N} P(B \cap A_i) = \sum_{i=1}^{N} P(A_i) P(B \mid A_i)$$
  \end{block}

  \vspace{0.2em}
  \pause
  \begin{alertblock}{Preparación para el Teorema de Bayes}
    Esta ley nos permite calcular la probabilidad del efecto final $B$ ponderando las probabilidades condicionales de cada causa posible $A_i$ por la probabilidad prior de dicha causa.
  \end{alertblock}
\end{frame}

% Slide 9: Independencia Estadística vs. Mutuamente Excluyentes
\begin{frame}{Independencia Estadística de Eventos}
  \small Si la ocurrencia de $A$ no aporta ninguna información ni modifica la probabilidad de $B$, es decir, si $P(B \mid A) = P(B)$, diremos que los eventos son \emph{independientes}. \pause

  \vspace{0.2em}
  \begin{block}{Definición de Independencia Estadística}
    Dos eventos $A, B \subset S$ son independientes si y solo si:
    $$P(A \cap B) = P(A) P(B)$$
  \end{block}

  \vspace{0.1em}
  \pause
  \begin{alertblock}{¡Precaución! Independiente $\neq$ Mutuamente Excluyente}
    \begin{itemize}\itemsep=0.1em
      \item \textbf{Mutuamente Excluyentes ($A \cap B = \emptyset$):} Si $A$ ocurre, $B$ \emph{no puede} ocurrir ($P(B \mid A) = 0$). ¡Son altamente dependientes!
      \item \textbf{Independientes:} Saber que $A$ ocurrió deja a $P(B)$ intacta.
    \end{itemize}
  \end{alertblock}
\end{frame}

% Slide 10: Independencia Dos a Dos vs. Independencia Mutua
\begin{frame}{Independencia Dos a Dos vs. Independencia Mutua}
  \footnotesize
  La independencia se generaliza a más de dos eventos, pero requiere extrema precisión en sus condiciones algebraicas: \pause

  \vspace{0.1em}
  \begin{block}{Independencia Mutua de Tres Eventos ($A_1, A_2, A_3$)}
    Tres eventos son mutuamente independientes si satisfacen simultáneamente \textbf{todas} las siguientes identidades de producto:
    \begin{align*}
      P(A_1 \cap A_2) &= P(A_1)P(A_2), \quad P(A_1 \cap A_3) = P(A_1)P(A_3) \\
      P(A_2 \cap A_3) &= P(A_2)P(A_3), \quad P(A_1 \cap A_2 \cap A_3) = P(A_1)P(A_2)P(A_3)
    \end{align*}
  \end{block}

  \vspace{0.1em}
  \pause
  \begin{block}{Observación}
    Es posible que tres eventos sean independientes por parejas (\emph{dos a dos}), pero que la ocurrencia conjunta de los dos primeros determine por completo al tercero ($P(A_3 \mid A_1 \cap A_2) = 1$). Por ello, la identidad del triple producto es indispensable.
  \end{block}
\end{frame}

% Slide 11A-1: Laboratorio Python (Parte 1: Urnas sin Reemplazo - Setup)
\begin{frame}[fragile]{Laboratorio en Python: Urna sin Reemplazo (1/2)}
  Simulación Monte Carlo de dos extracciones sucesivas sin reemplazo (\texttt{02.04\_probabilidad\_condicional.py}):
  \vspace{0.2em}
  {\lstset{basicstyle=\fontsize{6.3pt}{7.3pt}\selectfont\ttfamily}
  \lstinputlisting[firstline=23, lastline=39]{../../code/02_teoria_probabilidad/02.04_conditional_probability.py}}
\end{frame}

% Slide 11A-2: Laboratorio Python (Parte 1: Urnas sin Reemplazo - Lógica)
\begin{frame}[fragile]{Laboratorio en Python: Urna sin Reemplazo (2/2)}
  Conteo de casos condicionales e impresión de probabilidades empíricas vs. teóricas:
  \vspace{0.2em}
  {\lstset{basicstyle=\fontsize{6.3pt}{7.3pt}\selectfont\ttfamily}
  \lstinputlisting[firstline=40, lastline=56]{../../code/02_teoria_probabilidad/02.04_conditional_probability.py}}
\end{frame}

% Slide 11B: Laboratorio Python (Parte 2: Independencia en Dados)
\begin{frame}[fragile]{Laboratorio en Python: Test de Independencia}
  Lanzamiento de dos dados verificando empíricamente la dependencia en eventos acoplados por suma:
  \vspace{0.2em}
  {\lstset{basicstyle=\fontsize{6.3pt}{7.3pt}\selectfont\ttfamily}
  \lstinputlisting[firstline=75, lastline=94]{../../code/02_teoria_probabilidad/02.04_conditional_probability.py}}
\end{frame}

% Slide 11C: Laboratorio Python (Parte 3: Paradoja de Monty Hall)
\begin{frame}[fragile]{Laboratorio en Python: Paradoja de Monty Hall}
  Simulación de las estrategias del concurso demostrando el beneficio de cambiar de puerta:
  \vspace{0.2em}
  {\lstset{basicstyle=\fontsize{6.3pt}{7.3pt}\selectfont\ttfamily}
  \lstinputlisting[firstline=119, lastline=135]{../../code/02_teoria_probabilidad/02.04_conditional_probability.py}}
\end{frame}

% Slide 11D: Laboratorio Python (Salida en Consola)
\begin{frame}[fragile]{Laboratorio en Python: Salida en Terminal}
  Resultado exacto ejecutado en la consola de Python ($N = 100,000$):
  \vspace{0.2em}
  \begin{block}{Salida Estándar en Terminal}
    \fontsize{6.0pt}{7.0pt}\selectfont
    \begin{verbatim}
--- 1. Extracciones sin Reemplazo (Urna 5 Rojas, 3 Azules) ---
P(1a Roja) [A]:               Empirica = 0.6279 | Teorica = 0.6250 (5/8)
P(2a Roja | 1a Roja) [B|A]:   Empirica = 0.5729 | Teorica = 0.5714 (4/7)
P(2a Roja | 1a Azul) [B|A']:  Empirica = 0.7145 | Teorica = 0.7143 (5/7)
P(2a Roja Total) [B]:         Empirica = 0.6256 | Teorica = 0.6250

--- 2. Test de Independencia y Dependencia (Dos Dados) ---
P(Suma >= 9) [A]:             Empirica = 0.2797 | Teorica = 0.2778 (10/36)
P(Primer dado par) [B]:       Empirica = 0.5010 | Teorica = 0.5000 (18/36)
P(A inter B):                 Empirica = 0.1681 | Teorica = 0.1667 (6/36)
P(A | B):                     Empirica = 0.3355 | Teorica = 0.3333 (1/3)
Producto P(A) * P(B):         0.1401 != P(A inter B) -> ¡DEPENDIENTES!

--- 3. Paradoja de Monty Hall (100,000 partidas) ---
P(Ganar quedandose):          0.3326 (Teorica: 1/3 = 0.3333)
P(Ganar cambiando de puerta): 0.6674 (Teorica: 2/3 = 0.6667)
    \end{verbatim}
  \end{block}
\end{frame}

% Slide 12: Ejercicios del Libro — Nivel 1: Fundamental
\begin{frame}{Ejercicios del Libro — Nivel 1: Fundamental}
  \footnotesize
  \begin{block}{Problema 2.4.2 (Extracción Sucesiva sin Reemplazo)}
    Una urna contiene 5 bolas rojas y 3 azules. Se extraen dos bolas sin reemplazo. Sea $A$: \texttt{``la 1a es roja''} y $B$: \texttt{``la 2a es roja''}. Calcula $P(A)$, $P(B \mid A)$, $P(B \mid A')$ y la probabilidad total $P(B)$.
  \end{block}

  \vspace{0.1em}
  \pause
  \begin{alertblock}{Solución Analítica}
    \begin{itemize}\itemsep=0.1em
      \item $P(A) = \dfrac{5}{8} = 0.6250$. Si $A$ ocurre, quedan 4 rojas de 7: $P(B \mid A) = \dfrac{4}{7} \approx 0.5714$.
      \item Si $A'$ ocurre (1a azul), quedan las 5 rojas intactas de 7: $P(B \mid A') = \dfrac{5}{7} \approx 0.7143$.
      \item Por la Ley de Probabilidad Total:
      $$P(B) = P(A)P(B \mid A) + P(A')P(B \mid A') = \left(\frac{5}{8}\right)\left(\frac{4}{7}\right) + \left(\frac{3}{8}\right)\left(\frac{5}{7}\right) = \frac{20 + 15}{56} = \frac{5}{8}$$
    \end{itemize}
  \end{alertblock}
\end{frame}

% Slide 13: Ejercicios del Libro — Nivel 2: Operativo
\begin{frame}{Ejercicios del Libro — Nivel 2: Operativo}
  \footnotesize
  \begin{block}{Problema 2.4.4 (Control de Calidad en Líneas de Automatización)}
    El 70\% de las piezas provienen de la línea A y el 30\% de la línea B. La línea A tiene 5\% de defectos y la línea B un 10\%. Calcula la probabilidad total de que una pieza seleccionada al azar sea defectuosa ($D$).
  \end{block}

  \vspace{0.1em}
  \pause
  \begin{alertblock}{Desarrollo por Probabilidad Total}
    Identificamos la partición del espacio muestral como $S = A \sqcup B$, y los datos condicionales dados por el problema:
    $$P(A) = 0.70, \; P(B) = 0.30, \quad P(D \mid A) = 0.05, \; P(D \mid B) = 0.10$$
    Aplicando el Teorema 2.4.2:
    \begin{align*}
      P(D) &= P(A)P(D \mid A) + P(B)P(D \mid B) \\
           &= (0.70)(0.05) + (0.30)(0.10) = 0.035 + 0.030 = 0.065 \; (6.5\%)
    \end{align*}
  \end{alertblock}
\end{frame}

% Slide 14: Ejercicios del Libro — Nivel 3: Analítico
\begin{frame}{Ejercicios del Libro — Nivel 3: Analítico}
  \footnotesize
  \begin{block}{Problema 2.4.8 (Independencia en Complementos)}
    Demuestra formalmente que si los eventos $A$ y $B$ son estadísticamente independientes y $0 < P(A) < 1$, entonces sus complementos $A'$ y $B'$ también son mutuamente independientes.
  \end{block}

  \vspace{0.1em}
  \pause
  \begin{alertblock}{Demostración por Álgebra de Complementos}
    Queremos probar que $P(A' \cap B') = P(A')P(B')$. Por las Leyes de De Morgan, sabemos que $A' \cap B' = (A \cup B)'$. Por tanto:
    \begin{align*}
      P(A' \cap B') &= 1 - P(A \cup B) = 1 - [P(A) + P(B) - P(A \cap B)] \\
                    &= 1 - P(A) - P(B) + P(A)P(B) \quad \text{(por hipótesis de indep.)} \\
                    &= [1 - P(A)] - P(B)[1 - P(A)] = [1 - P(A)][1 - P(B)] \\
                    &= P(A')P(B') \quad \blacksquare
    \end{align*}
  \end{alertblock}
\end{frame}

% Slide 15: Ejercicios del Libro — Nivel 4: Desafiante (1/2)
\begin{frame}{Ejercicios del Libro — Nivel 4: Desafiante (1/2)}
  \footnotesize
  \begin{block}{Problema 2.4.10: La Paradoja de Monty Hall (Enunciado)}
    En un concurso hay 3 puertas ($P_1, P_2, P_3$): 1 oculta un auto y 2 ocultas cabras. El participante elige la Puerta 1. El presentador, quien sabe dónde está el auto y siempre debe abrir una puerta con cabra, abre la Puerta 3. ¿Debe el participante cambiar su elección a la Puerta 2?
  \end{block}

  \vspace{0.1em}
  \pause
  \begin{alertblock}{Formalización Axiomática de la Hipótesis Previa}
    Sea $C_i$: \texttt{``El auto está tras la puerta $i$''} ($P(C_1)=P(C_2)=P(C_3)=1/3$).
    Sea $A_3$: \texttt{``El presentador abre la Puerta 3''}. Evaluamos las condicionales:
    \begin{itemize}\itemsep=0.05em
      \item Si $C_1$ (auto en 1), el presentador elige al azar entre 2 y 3: $P(A_3 \mid C_1) = 1/2$.
      \item Si $C_2$ (auto en 2), el presentador \emph{está obligado} a abrir la 3: $P(A_3 \mid C_2) = 1$.
      \item Si $C_3$ (auto en 3), el presentador no puede abrir la 3: $P(A_3 \mid C_3) = 0$.
    \end{itemize}
  \end{alertblock}
\end{frame}

% Slide 16: Ejercicios del Libro — Nivel 4: Desafiante (2/2)
\begin{frame}{Ejercicios del Libro — Nivel 4: Desafiante (2/2)}
  \footnotesize
  Calculamos primero la probabilidad total de que el presentador abra la Puerta 3: \pause
  $$P(A_3) = P(C_1)P(A_3 \mid C_1) + P(C_2)P(A_3 \mid C_2) + P(C_3)P(A_3 \mid C_3) = \frac{1}{3}\left(\frac{1}{2}\right) + \frac{1}{3}(1) + 0 = \frac{1}{2}$$ \pause

  \vspace{0.1em}
  \begin{block}{Comparación Exacta de Estrategias}
    \begin{itemize}\itemsep=0.1em
      \item \textbf{Probabilidad si se queda en Puerta 1:}
      $$P(C_1 \mid A_3) = \dfrac{P(C_1 \cap A_3)}{P(A_3)} = \dfrac{P(C_1)P(A_3 \mid C_1)}{P(A_3)} = \dfrac{(1/3)(1/2)}{1/2} = \mathbf{\dfrac{1}{3}}$$
      \item \textbf{Probabilidad si cambia a la Puerta 2:}
      $$P(C_2 \mid A_3) = \dfrac{P(C_2 \cap A_3)}{P(A_3)} = \dfrac{P(C_2)P(A_3 \mid C_2)}{P(A_3)} = \dfrac{(1/3)(1)}{1/2} = \mathbf{\dfrac{2}{3}}$$
    \end{itemize}
    ¡Cambiar de puerta duplica matemáticamente las probabilidades de ganar!
  \end{block}
\end{frame}

% Slide 17: Síntesis y Reflexión
\begin{frame}{Cierre de Sección y Síntesis Condicional}
  \begin{reflexion}[El Poder de la Actualización Dinámica]
    La probabilidad condicional no es simplemente una fórmula de división; es el mecanismo formal mediante el cual el conocimiento riguroso actualiza nuestras creencias. Transformar el universo muestral $S$ en el evento confirmado $A$ elimina la incertidumbre estática y da paso al modelado inteligente de datos reales.
  \end{reflexion}

  \vspace{0.5em}
  \pause
  \begin{retoclase}[Para la próxima sesión: Teorema de Bayes]
    En esta sección calculamos el efecto a partir de las causas mediante la Ley de la Probabilidad Total. En la próxima sesión resolveremos el problema inverso del diagnóstico científico: dada la observación del efecto ($B$), ¿cuál es la probabilidad exacta de cada una de sus posibles causas ($A_k$)?
  \end{retoclase}
\end{frame}

\end{document}
